% 09 · Validation — how to know whether any of this is real
\section{VALIDATION: KNOWING WHEN A SCORE LIES}
\label{sec:validation}

\deckslides{slides 19--20 (internal and external validation) and slide 3.}

In a supervised problem, validation is arithmetic: hold out labels, count
mistakes. In clustering there is nothing to hold out, and the score you report
is chosen by you. This chapter is therefore the most important one in the
workbook, and the one in which we report the most unflattering numbers ---
including two of our own that forced us to withdraw a published claim.

\subsection{Internal indices: is the partition tidy?}

Three classical measures, all computed from the data and the partition alone.

\begin{itemize}
\item \textbf{Sum of squared errors.} The K-means objective of
  \textbf{Equation~\ref{eq:sse}} (\S\ref{sec:kmeans}).
  Monotone in the number of groups, so it is only useful as a curve against
  $K$, never as a score.
\item \textbf{Silhouette} \citep{Rousseeuw:87}. For a point $i$, let $a(i)$ be
  its mean distance to the other members of its own group and $b(i)$ its mean
  distance to the members of the nearest other group. Then
  $s(i) = (b(i)-a(i))/\max\{a(i),b(i)\} \in [-1,1]$, and the silhouette is the
  mean over points. It rewards compact, well-separated groups --- which is a
  statement about geometry, not about truth.
\item \textbf{Dip test} \citep{Hartigan:85}. A significance test for
  unimodality of a one-dimensional distribution. Its role here is upstream of
  clustering: it answers whether a variable (or a projection) shows any
  multipeaked structure at all. If nothing is multimodal, the clusters an
  algorithm returns are imposed by the algorithm.
\end{itemize}

All three share one weakness, and it is fatal if forgotten: \emph{a tidy
partition of unstructured data is easy to produce}. The silhouette of a
$K$-means solution on uniform noise is not zero, and the dip test is one
significance test among many you could run. Internal indices are for
diagnosing the behaviour of an algorithm, not for certifying a discovery.

\begin{marginnote}[]
\entry{Homogeneity}{How pure each predicted group is: 1 when every cluster an
algorithm reports contains a single true cluster, even if that true cluster is
split across several predictions.}
\entry{Completeness}{How intact each true cluster is: 1 when every true
cluster ends up in one predicted group, even if that group contains several
true clusters. Homogeneity and completeness trade off.}
\end{marginnote}

\subsection{External indices: comparison with something the algorithm never saw}

Here the labels come from outside the algorithm. Two families are used in this
workbook, and we use them for different tasks.

\textbf{Cluster-only separation} follows \citet{GarciaDias:19}. Take only the
known members, cluster them, and compare the predicted grouping with the true
one using an information-theoretic pair plus an assignment score:

\begin{align}
h &= 1 - \frac{H(C \mid K)}{H(C)}, &
c &= 1 - \frac{H(K \mid C)}{H(K)}, &
v &= \frac{2hc}{h+c},
\label{eq:merits}
\end{align}

where $C$ are the true cluster labels, $K$ the predicted ones, and $H$ the
Shannon entropy; accuracy is the fraction of correctly assigned stars after
matching predicted to true labels with the Hungarian algorithm (maximising the
diagonal of the confusion matrix). The pairing is not cosmetic. A method that
merges every cluster into one group has $c = 1$ and $h \approx 0$; a method
that shatters them has $h = 1$ and low $c$. Quoting one without the other is
how a broken result reads as a good one, and we do it ourselves in
\textbf{Table~\ref{tab:honest}} for exactly that reason.

\textbf{Field retrieval} follows the operational question: for each cluster,
take the predicted group that overlaps its true members most, and measure

\begin{equation}
\text{recall} = \frac{\text{true members recovered}}{\text{all true members}},
\qquad
\text{precision} = \frac{\text{true members in the group}}{\text{stars in the group}} .
\label{eq:rp}
\end{equation}

\textbf{Recovery fraction} \citep{Casamiquela:21} answers the question a
per-cluster study actually asks: how many of the known clusters came back? For
each true cluster, take the predicted group that holds most of its stars (the
best-overlap rule above). The cluster counts as \emph{recovered} at a threshold
$\theta$ when that group holds at least $\theta$ of the cluster \emph{and} at
least $\theta$ of the group is the cluster --- their per-cluster completeness and
homogeneity, our recall and precision, both at once:

\begin{equation}
\mathrm{RF}_{\theta} = \frac{\#\{\text{clusters recovered at } \theta\}}{\#\{\text{clusters}\}},
\qquad \theta = 0.40 \text{ or } 0.70 .
\label{eq:rf}
\end{equation}

Its companion is the \emph{statistical fraction}: the share of predicted groups
that recover no cluster at all, which is how the other paper measures the amount
of decoration in its own solution. We adopt this pair for the same reason we
adopt their task: of all the scores in this chapter it is the one whose chance
level is near zero whatever the sizes and number of clusters, so it survives
comparison between different samples. The implementation is
\texttt{cluster.baseline.recovery\_fraction}, and it reproduces their Table~2
exactly, which is the only external check on it available.

Two complements to \textbf{Equation~\ref{eq:rp}} appear in the results.
\emph{kNN purity} (\S\ref{sec:knn}) is parameter-free cohesion, immune to the
best-overlap protocol below. \emph{AUROC} answers a refusal question rather
than a discovery one: given a true member and a contaminant, how often does
the method rank the member as more plausible? It is used in
\S\ref{sec:benchmark} to show that the rejection power of spectra grows with
cluster age.

\begin{issues}[THE ``BEST OVERLAP'' IS A BEST CASE]
The protocol above scores a method by its single best-matching predicted
group per cluster. That is deliberately generous --- it is the only fair way to
compare algorithms that return different numbers of groups --- but it means a
method can look good by returning many groups, only one of which overlaps each
cluster. It also cannot be gamed into positive precision: the denominator is
the size of the predicted group, so a group with 3 true members in it and 300
field stars scores 0.01 however well it ``recovered'' those three.
\end{issues}

\subsection{Five ways a score lies}

Each of these is a case we hit. They are ordered from the oldest to the most
recent, and the fourth is the reason the validation chapter is in the middle
of the workbook rather than at the end.

\paragraph{1 · The degenerate partition.} A clustering that returns two groups
with two thirds of the stars in one of them produces a homogeneity of
$0.269$ --- a floor set by the collapse, not a measurement of the method. Two
of our PCA rows sat on that floor, and a first draft of this work described
the spectral latent as ``three times better than PCA''. It was three times
better than a crash. The check that catches it is a column nobody prints by
default: the fraction of stars in the largest predicted group. We now print
it, and \texttt{cluster head-to-head} flags degenerate rows automatically.

\paragraph{2 · Comparing different populations.} The abundance arm of our
first head-to-head was scored on 25 clusters and 1\,215 stars, while every
other arm was scored on 5 clusters and 55 stars. The abundance row therefore
looked worse than the others for a reason that had nothing to do with
abundances: it was being asked a harder question. The fix is structural, not
statistical --- intersect the sets of stars before scoring, so every row of
the table describes the same 982 stars and the same 25 clusters. Two tables in
this workbook have been rebuilt that way, and the rebuilt numbers are the ones
quoted.

\paragraph{3 · Row order, and duplicate stars.} Two checks on the same member
matrix, both of which should have been run before any of it was published.

\emph{Row order.} Permuting the rows of the abundance matrix --- a
transformation that leaves every pairwise distance exactly unchanged --- moves
the t-SNE homogeneity between $0.218$ and $0.560$, and the number of predicted
groups between 2 and 22. The sorted order gives $0.560$; the mean over eight
random orders is $0.268 \pm 0.100$. The geometry is identical in all nine
runs; what changes is the optimiser's path, which depends on initialisation
and on the order in which the Barnes--Hut approximation accumulates
interactions. Every other arm is stable to a few hundredths: UMAP on the same
abundances $0.579 \pm 0.009$, and the masked-AE latent $0.740 \pm 0.003$ under
t-SNE. The lesson is not ``t-SNE is broken'' but ``report the row-order spread
of any embedding with a stochastic initialisation'', because on small samples
it can exceed the effect you are claiming.

\emph{Duplicates.} An audit of the same matrix found that it contains repeat
entries for the same star: of 1\,002 member rows, 807 carry distinct
\texttt{APOGEE\_ID}s and 34 carry none, so 332 rows share an identifier with
another row (\texttt{scripts/population\_counts.py}). The copies are not identical --- the median spread between two
rows of the same star is $0.04$ dex per element, comparable to the abundance
differences we are trying to measure --- and they are concentrated: M~3
contributes 102 duplicate rows, M~67 ninety. This inflates the member counts
of \textbf{Table~\ref{tab:clusters}} (they are row counts, not star counts),
double-weights those clusters in every macro average, and is the most likely
explanation for the fragility of the abundances-only t-SNE row above. The
clean fix is to collapse each star's rows before clustering.

\paragraph{4 · Tuning on the score.} The lever sweep of
\S\ref{sec:benchmark} chose element weights by maximising the very metric it
then reported, on a single cluster at a single seed. That is a legitimate way
to \emph{find} a lever and an illegitimate way to \emph{quantify} it, and the
current run exposed the difference: the same switch that lifted t-SNE recall
on the DR17 data ($0.09 \rightarrow 0.44$) lowers precision on the DR19 data
($0.83 \rightarrow 0.24$). A lever is a hypothesis; only a re-run on the
current data, quoted with both recall and precision, turns it into a result.

\paragraph{5 \textperiodcentered\ The chance level of a mutual-information score moves
with the sample.} Small groups inflate $h$, $c$ and $V$. Shuffle
\citet{Casamiquela:21}'s 175 stars at random into groups with the sizes they
observe --- 31 clusters, most of them tiny --- and the V-measure is $0.51$; do the
same on our 982 stars in 25 clusters and it is $0.06$--$0.12$. Their published
$V = 0.55$ therefore sits $0.04$ above its null, and our $V = 0.58$ sits about
$0.47$ above ours. The two numbers are not comparable, and a table that prints
them side by side without their nulls asserts a comparison that does not exist.
\textbf{Equation~\ref{eq:rf}} does not have this problem: chance level $0.04$
for their sample shape and $\le 0.001$ for ours. This is the reason
\S\ref{sec:literature} is written in recovery fractions.

\subsection{The protocol this workbook follows}

Out of all that, six rules, which the numbers in
\S\ref{sec:benchmark}--\S\ref{sec:physics} obey:

\begin{enumerate}
\item \textbf{Same population or no comparison.} Arms are intersected on the
  stars they share before any score is computed.
\item \textbf{Mean $\pm$ standard deviation over seeds.} Seven seeds
  ($42, 0, 1, 2, 7, 13, 99$), never a single run, and never a best-of.
\item \textbf{Kinematics are the ceiling, not a competitor.} They define the
  labels; the interesting question is always how close chemistry or spectra
  get to that ceiling.
\item \textbf{Parameter-free scores alongside tuned ones.} kNN purity and
  AUROC cannot be improved by tuning, which makes them the sanity check on
  everything else.
\item \textbf{Report the degenerate rows.} Largest-group fraction and number
  of predicted groups are printed next to every score, so a collapse is
  visible without recomputation.
\item \textbf{Recompute the chance level before comparing with a published
  score.} A score quoted without its null is a number without units, and the
  null moves when the sample does (rule 5 above).
\end{enumerate}

\begin{figure}[htbp]\centering
  \begin{tabular}{@{}c@{\hskip8pt}c@{}}
    \fig[0.44]{confusion_tsne_chem.png} &
    \fig[0.44]{confusion_tsne_kin.png}
  \end{tabular}
  \caption{The confusion matrix of a t-SNE clustering of the members, with the
  true cluster on one axis and the predicted group on the other
  (\textbf{Equation~\ref{eq:merits}} visualised). Left: abundances only --- one
  predicted group absorbs most of the true clusters, which is the
  completeness-rich, homogeneity-poor blob described above. Right: adding
  kinematics. The blocks along the diagonal are what ``the clusters are
  separated'' looks like in a matrix, and the difference between the two panels
  is the central result of \S\ref{sec:benchmark}.}
  \label{fig:confusion}
\end{figure}

\begin{figure}[htbp]\centering
  \fig[0.78]{proper_motions.png}
  \caption{Why kinematics can referee. Each panel shows proper motion in right
  ascension and declination for the stars around one cluster; the tight group
  sharing a velocity is the cluster, the surrounding cloud is the field. These
  are the observables that define the labels, so they can never be an input to
  a method that is scored against them.}
  \label{fig:propermotions}
\end{figure}

\begin{issues}[EXERCISES]
\begin{enumerate}
\item Show that a clustering that puts every star in a single group has
  $h = 0$ and $c = 1$, and that a clustering that gives every star its own
  group has $h = 1$ and $c \approx 0$. Where does accuracy sit in each limit?
\item Reproduce the row-order check: cluster the DR19 member matrix with
  t-SNE$\,+\,$HDBSCAN nine times, once sorted and eight times with a random
  row order, fixing the seed. Report the spread of homogeneity and of the
  number of predicted groups. Now do the same with the masked-AE latent. Why
  is one stable and the other not?
\item Collapse the duplicate rows of the member matrix (for example by taking
  the median abundance vector per \texttt{APOGEE\_ID}) and re-run the
  cluster-only table. Which rows of \textbf{Table~\ref{tab:honest}} move, and
  does M~3's share of the sample change the conclusion of the ablation?
\item Design a $\sigma$-clipped chemical membership and use it to score a
  kinematic clustering. Based on \S\ref{sec:data}, explain in three sentences
  why the resulting number would be meaningless --- and what the equivalent
  mistake would look like in your own field.
\end{enumerate}
\end{issues}
